%%
%% Author: shoupu_wan
%% 7/14/18
%%
\chapter{Reusable Functions}

For simplicity, I refer to standalone functions as well as member methods of a class as
\emph{functions}.

\section{Create easy to reuse modules or functions}

\label{Evernote:Rotate array by k times}
\label{Evernote:Search for start and end of a target in a sorted array}

\section{Why function?}

It is the same answer we would give to the question "Why design?", that is, to support abstract
coding.
Our brains are small.
So we have to reduce the complexity of code to understand the symantics.
Abstract the implementation with the name of a function enables us to do works with more complexity.
More importantly, think of coding as building pyramids.
Firstly we make the building blocks-bricks, carve stones, etc..
Then we build structures with these blocks.
Finally, we assemble structures into pyramid.
Higher level constructions depends on lower level constructions.

\section{Case Study--Spiral Matrix} \label{sec:case-spiral-matrix}

\label{Evernote:Refactor process example - Spiral Matrix}

\section{Spiral printing of matrix elements}

Spiral Printing of an 2D matrix. See EverNote

\section{Modularity}

Just breaking up a 100-line of function into five 20-line function does not create modularity.
That only makes it worse.

Here are some tips about you should break out a chunk of code into a separate function
\begin{itemize}
    \item In retrospective, understanding the code you just typed is beyond a glimpse.
    \item You are embedding a subtask $B$ within the body of your main function
    $A$.
    \item Obtaining certain information is so lengthy that it obstructs the main flow of control.
    \item A squad of certain lines are needed in more than two places, near or far, in your codebase.
    \item A sequence of operations appears again and once more.
    \item You encounter difficulty to create effective tests for a function.
\end{itemize}

These are low-hanging fruits of signals